\chapter{Schach}

\section{Einstufung der Spieler nach Elo-Zahl}

Vor Einführung der Elo-Zahl stufte man die Spieler beim Schach in neun Klassen oder Kategorien ein. Ein Unterschied von einer Klasse bedeutete, dass der bessere Spieler als Ergebnis einer Partie $0,75$ Punkte \textit{erwarten} darf. Im Elo-System entspricht dieser Spielstärkeunterschied einer Differenz von etwa $200$ Wertungspunkten (genauer: $191$, s. Anm. \ref{anm_5}). Eine Erweiterung stellt das Glicko-System dar.

\begin{table}[htbp]
    \centering
    \begin{tabular}{c >{\centering\arraybackslash}p{6.2cm} >{\centering\arraybackslash}p{6.2cm}}
        \toprule
        \textbf{Elo-Zahl} & \textbf{Offene Kategorie} & \textbf{Kategorie Frauen} \\
        \midrule
        $\geq 2500$     & \multicolumn{2}{c}{Großmeister (GM)} \\
        $2400 - 2499$  & \multicolumn{2}{c}{Internationaler Meister (IM)} \\
        $2300 - 2399$  & FIDE-Meister (FM) & Großmeister der Frauen (WGM) \\
        $2200 - 2299$  & Candidate Master (CM) oder Nationaler Meister & Internationaler Meister der Frauen (WIM) \\
        $2100 - 2199$  & Meisteranwärter & FIDE-Meister der Frauen (WFM) \\
        $2000 - 2099$  & Starker Amateur & Candidate Master der Frauen (WCM) \\
        $< 2000$        &  \multicolumn{2}{c}{Amateur} \\
        \bottomrule
    \end{tabular}
    \caption{Zuordnung der Titel nach der Wertungszahl.}
    \label{tab:schach-titel}
\end{table}

Zu beachten ist dabei, dass man die verschiedenen Titel Großmeister (GM) und Internationaler Meister (IM) nicht nur auf Grund einer bestimmten Elo-Zahl erhält, sondern durch die Erfüllung von anderen festgelegten Normen. Um den Titel nach Erfüllung aller Normen zu erhalten, muss ein angehender GM allerdings eine Elo-Zahl von mindestens $2500$, ein IM eine Zahl von mindestens $2400$ einmal erreicht haben. Bei entsprechender Punktzahl wird der offene Großmeistertitel an Spieler aller Geschlechter vergeben, der Großmeistertitel der Frauen kann von diesen bereits ca. $200$ Punkte vorher erreicht werden (\autoref{tab:schach-titel}).

Unterhalb $2000$ Elo-Punkten gibt es keine offiziellen Titel, weshalb eine objektive Benennung der Spielstärke schwierig ist. Zur groben Einschätzung kann in Deutschland der DWZ-Schnitt der Amateurligen dienen (Die Deutsche Wertzahl DWZ ist zwar nicht dasselbe wie die international übliche Elo-Zahl, aber vergleichbar). Im Schachverband Württemberg liegt der DWZ-Schnitt in der Oberliga bei ca. $2100$, in den Verbandsligen bei ca. $1900$, in den Landesligen bei ca. $1700$, den Bezirksligen bei ca. $1600$ usw. Wie in anderen Sportarten auch erreichen Gelegenheitsspieler selten das Niveau der Bezirksliga.

\begin{wraptable}{r}{0.37\textwidth}
    \centering
    \begin{tabular}{c c c}
        \toprule
        \textbf{Turnier-} & \multicolumn{2}{c}{\textbf{Elo-Durchschnitt}} \\
        \textbf{kategorie} & \textbf{Von} & \textbf{Bis} \\
        \midrule
        \textbf{1} & $2251$ & $2275$ \\
        \textbf{2} & $2276$ & $2300$ \\
        \textbf{3} & $2301$ & $2325$ \\
        \multicolumn{3}{c}{$\dots$} \\
        \textbf{7} & $2401$ & $2425$ \\
        \multicolumn{3}{c}{$\dots$} \\
        \textbf{11} & $2501$ & $2525$ \\
        \multicolumn{3}{c}{$\dots$} \\
        \textbf{15} & $2601$ & $2625$ \\
        \multicolumn{3}{c}{$\dots$} \\
        \textbf{19} & $2701$ & $2725$ \\
        \textbf{20} & $2726$ & $2750$ \\
        \textbf{21} & $2751$ & $2775$ \\
        \textbf{22} & $2776$ & $2800$ \\
        \textbf{23} & $2801$ & $2825$ \\
        \bottomrule
    \end{tabular}
    \caption{Turnier-Kategorien.}
    \label{tab:turnier-kategorien}
\end{wraptable}

Nachdem im Jahr 1970 die Elo-Zahl als Wertungssystem eingeführt worden war, hatte zunächst Bobby Fischers Bestmarke von $2785$ Punkten vom Juli 1972 für viele Jahre Bestand. Im Jahre 1999 erreichte der damalige klassische Schachweltmeister Garri Kasparow die Elo-Zahl von $2851$ Punkten, die erst im Januar 2013 von Magnus Carlsen mit $2861$ Punkten übertroffen wurde. Inzwischen konnte Carlsen den Rekord auf $2882$ erhöhen (Mai 2014 und August 2019).

Spieler mit einer Elo-Zahl über $2600$ gehören zum erweiterten Kreis der Weltspitze (etwa TOP $200$), Spieler mit über $2700$ zum engeren Kreis (etwa TOP $30$). Aktuell (Stand: Dezember 2025) erreichen nur zwei Spieler $2800$, und zwar Hikaru Nakamura sowie Magnus Carlsen.

Elo-Zahlen können auch für einzelne Turniere berechnet werden (s.\,o.: \textit{Elo-Performance}). Hierbei gelang Fabiano Caruana im Jahr 2014 beim Turnier um den Sinquefield Cup in St. Louis eine Elo-Leistung von $3103$. Die davor gültige höchste Turnier-Elo-Leistung war $3002$, erzielt von Magnus Carlsen in Nanjing im Jahr 2010. \cite{loeffler2014}


\section{Historische Elo-Zahl im Schach}

Für den Vergleich heutiger Spitzenspieler mit Großmeistern vor der Einführung der Elo-Zahl wird die sogenannte historische Elo-Zahl verwendet. %TODO: ggf Hauptartikel verlinken

\section{Spielstärke der Schachprogramme}

Die Elo-Zahlen der Schachcomputer bzw. Computerprogramme sind nicht ohne weiteres mit denen menschlicher Schachspieler zu vergleichen, da sie überwiegend durch Partien zwischen Computern ermittelt wurden und nicht durch Teilnahme an offiziellen Turnieren. %TODO: ggf Hauptartikel verlinken


\section{Turnierkategorien -- Einteilung der Turniere nach Elo-Zahl}

Auch Rundenturniere werden nach der durchschnittlichen Elo-Zahl der Teilnehmer in Kategorien eingeteilt. Hierbei entspricht ein Unterschied um eine Kategorie $25$ Elo-Punkten (\autoref{tab:turnier-kategorien}). Als Turnier der Kategorie 1 wird dabei ein Turnier eingestuft, dessen Teilnehmer im Durchschnitt $2251$ bis $2275$ Elo-Punkte haben. Die zurzeit stärksten Turniere erreichen die Kategorie 22, was einem Durchschnitt von $2776$ bis $2800$ Elo-Punkten entspricht. Bei der Zürich Chess Challenge 2014 wurde im Januar 2014 erstmals Kategorie 23 (mit einem Elo-Durchschnitt von $2801$) erreicht.